%==============================================================================
\chapter{Introducción}
\label{cap:introduccion}
%==============================================================================

%------------------------------------------------------------------------------
\section{Contexto y planteamiento del problema}
\label{sec:contexto}
%------------------------------------------------------------------------------
% Fuente: Anteproyecto, Cap. 1 (Introducción) y Cap. 2 (Planteamiento).

El Job Shop Scheduling Problem (JSSP) es uno de los problemas más estudiados en
el campo de la optimización combinatoria. En este problema, un conjunto de
trabajos debe ser procesado en distintas máquinas, cumpliendo para cada trabajo
una secuencia específica de operaciones. El objetivo suele ser minimizar el
tiempo total de finalización (\textit{makespan}), aunque pueden considerarse
otros criterios como el retraso promedio o la utilización de recursos
\citep{pinedo2022}.

El JSSP es un problema NP-duro, por lo que su resolución eficiente requiere
algoritmos avanzados, entre ellos heurísticas, metaheurísticas y métodos
híbridos. Sin embargo, el rendimiento de estos algoritmos depende en gran
medida de las características de cada instancia, lo que hace difícil determinar
de antemano cuál ofrecerá los mejores resultados. En la práctica, seleccionar
manualmente el algoritmo más apropiado es ineficiente y requiere experiencia
especializada.

Para abordar esta incertidumbre, los portafolios de algoritmos se han
convertido en una estrategia efectiva, puesto que combinan múltiples métodos y
seleccionan aquel que se espera tenga el mejor desempeño
\citep{smithmiles2009}. No obstante, la mayoría de los enfoques adoptan la
política \textit{Winner Takes All}, que consiste en seleccionar un solo
algoritmo por instancia. Este enfoque es vulnerable a errores de selección y no
aprovecha la complementariedad entre algoritmos. Por ello, junto a selectores por
instancia como SATzilla \citep{xu2008satzilla}, han surgido métodos como
CPHydra \citep{omahony2008cphydra} y SUNNY \citep{amadini2014sunny}, que
emplean técnicas de aprendizaje automático para definir un \textit{schedule}
de ejecución, es decir, el orden y el tiempo asignado a cada algoritmo dentro
del portafolio. Estos sistemas han demostrado
mejorar el rendimiento global y la robustez en la resolución de problemas
complejos.

Este trabajo aplica esa idea al JSSP clásico mediante el diseño de un
portafolio de algoritmos y un esquema de ejecución planificado con
\textit{Machine Learning}, capaz de decidir dinámicamente la secuencia y la
duración de ejecución de los métodos más apropiados para cada instancia.

%------------------------------------------------------------------------------
\section{Formulación del problema}
\label{sec:formulacion}
%------------------------------------------------------------------------------
% Fuente: Anteproyecto, Cap. 2 (Pregunta problema).

La pregunta que guía el trabajo es la siguiente:

\begin{quote}
  ¿Cómo diseñar un \textit{schedule} de ejecución basado en portafolios que
  determine el orden y tiempo óptimo de ejecución de algoritmos, utilizando
  \textit{Machine Learning}, para mejorar la eficiencia y robustez en la
  resolución del Job Shop Scheduling Problem clásico?
\end{quote}

Para responderla, el trabajo construye un portafolio de algoritmos
complementado por un modelo de selección automática que indica en qué orden
ejecutar los algoritmos y durante cuánto tiempo, es decir, el \textit{schedule}
de ejecución.

%------------------------------------------------------------------------------
\section{Justificación}
\label{sec:justificacion}
%------------------------------------------------------------------------------
% Fuente: Anteproyecto, Cap. 3.1 (económica) y 3.2 (académica).

En el ámbito académico, el trabajo aporta al estudio del problema de selección
de algoritmos (\textit{Algorithm Selection Problem}), una línea de
investigación que combina el aprendizaje automático y la optimización
combinatoria. Ofrece un entorno experimental para explorar la caracterización
de instancias del JSSP, la aplicación de técnicas de clasificación y regresión
para predecir el rendimiento de los algoritmos, y la integración de modelos de
aprendizaje con métodos de optimización tradicionales, como la búsqueda tabú,
los algoritmos genéticos o el recocido simulado.

Desde la perspectiva económica y productiva, una planificación ineficiente en
entornos industriales genera costos por tiempos muertos, sobreutilización de
recursos y retrasos en la entrega de productos. Un portafolio inteligente puede
reducir los costos operativos mediante una mejor asignación de recursos,
disminuir el tiempo total de procesamiento (\textit{makespan}) e incrementar la
productividad de la planta. En el plano computacional, seleccionar
adecuadamente el algoritmo reduce tanto el tiempo de cómputo como el consumo de
recursos.

%------------------------------------------------------------------------------
\section{Objetivos}
\label{sec:objetivos}
%------------------------------------------------------------------------------
% Fuente: Anteproyecto, Cap. 4.1 y 4.2 (texto aprobado, no modificar sin
% autorización del director).

\subsection{Objetivo general}

Desarrollar un \textit{schedule} de ejecución basado en portafolios para el Job
Shop Scheduling Problem clásico, utilizando técnicas de \textit{Machine
Learning}.

\subsection{Objetivos específicos}

\begin{enumerate}[label=\arabic*., leftmargin=1.27cm, itemsep=0.3em]
  \item Analizar y comparar distintos enfoques de portafolios de algoritmos
    aplicados al Job Shop Scheduling Problem, incluyendo estrategias
    \textit{Winner Takes All} y planificación secuencial.
  \item Identificar y extraer características relevantes de las instancias del
    JSSP.
  \item Implementar un portafolio que integre la predicción y selección
    automática del algoritmo más adecuado.
  \item Validar experimentalmente el portafolio mediante pruebas con instancias
    de referencia.
\end{enumerate}

%------------------------------------------------------------------------------
\section{Alcance y limitaciones}
\label{sec:alcance}
%------------------------------------------------------------------------------
% Fuente: Anteproyecto, Cap. 5.

El trabajo se centra en el estudio y la resolución del Job Shop Scheduling
Problem clásico, con la minimización del \textit{makespan} como criterio
principal de optimización. Para ello se trabaja con un conjunto limitado pero
representativo de algoritmos, que abarca métodos heurísticos, metaheurísticos y
exactos, con el propósito de evaluar su comportamiento ante diferentes tipos de
instancias. El enfoque se basa en modelos de aprendizaje supervisado, que se
entrenan y validan con conjuntos de datos de referencia.

El portafolio se limita a la versión clásica del problema. No incluye variantes
como el Flexible Job Shop ni las versiones multiobjetivo, que se consideran
líneas de trabajo futuro.

En la práctica, el conjunto de algoritmos quedó compuesto por dos
\textit{solvers} exactos de programación por restricciones, Chuffed y Gecode,
en tres configuraciones, junto con una heurística de despacho que calcula una
cota superior del \textit{makespan} y una solución de respaldo. OR-Tools CP-SAT
se usa como \textit{solver} de referencia fuera del portafolio. La evaluación se
hace sobre 242 instancias clásicas de la literatura con un límite de 300
segundos por instancia. Las limitaciones que se derivan de estas decisiones se
discuten en la Sección~\ref{sec:limitaciones}.

%------------------------------------------------------------------------------
\section{Resultados obtenidos y contribuciones}
\label{sec:resultados_contribuciones}
%------------------------------------------------------------------------------
% Fuente: Anteproyecto, Cap. 4.3 y Tabla 4.1. Al cierre del trabajo, la
% columna de resultados esperados se reemplaza por los resultados obtenidos.

La Tabla~\ref{tab:objetivos_resultados} relaciona cada objetivo específico con
el resultado obtenido y con el capítulo de este documento donde se presenta.

\begin{table}[H]
  \centering
  \small
  \singlespacing
  \caption{Objetivos específicos, resultados obtenidos y capítulo donde se presentan.}
  \label{tab:objetivos_resultados}
  \begin{tabularx}{\textwidth}{>{\raggedright\arraybackslash}p{0.30\textwidth}
                               >{\raggedright\arraybackslash}X
                               >{\centering\arraybackslash}p{0.13\textwidth}}
    \toprule
    \textbf{Objetivo específico} & \textbf{Resultado obtenido} & \textbf{Capítulo} \\
    \midrule
    1. Analizar y comparar enfoques de portafolios (\textit{Winner Takes All} y
    planificación secuencial)
      & Comparación de \textit{Winner Takes All}, SUNNY, un \textit{schedule}
        estático y un \textit{schedule} con aprendizaje automático, en calidad
        de solución, óptimos probados, tiempo (PAR10) y robustez.
      & \ref{cap:portafolio}, \ref{cap:resultados} \\
    \addlinespace
    2. Identificar y extraer características de las instancias del JSSP
      & 65 características de \texttt{mzn2feat} y 20 propias del JSSP, junto
        con una base de conocimiento que incluye el desempeño y las
        trayectorias de cada configuración.
      & \ref{cap:caracterizacion} \\
    \addlinespace
    3. Implementar un portafolio con predicción y selección automática del
    algoritmo
      & Prototipo que calcula cotas, extrae características y construye el
        orden y el tiempo de ejecución por instancia con tres modelos de
        aprendizaje automático.
      & \ref{cap:portafolio} \\
    \addlinespace
    4. Validar experimentalmente el portafolio con instancias de referencia
      & Validación \textit{leave-one-out} sobre 236 instancias y tabla
        comparativa entre las configuraciones individuales con y sin cotas,
        los métodos de portafolio y el VBS.
      & \ref{cap:diseno}, \ref{cap:resultados} \\
    \bottomrule
  \end{tabularx}
\end{table}

Las principales contribuciones del trabajo son:

\begin{itemize}
  \item Un modelo con cotas para el JSSP en MiniZinc, que combina una cota
    inferior con la solución de la heurística de Giffler y Thompson y reduce el
    \textit{gap} medio de Chuffed de 648\,\% a 17\,\%. La comparación con los
    mismos \textit{solvers} sin cotas muestra que, sin ellas, ni siquiera
    eligiendo en cada instancia la mejor configuración el \textit{gap} medio
    baja del 24.5\,\%, mientras que con cotas todos los enfoques de portafolio
    quedan cerca del 13\,\%.
  \item Una base de conocimiento con trayectorias, que registra cuándo encontró
    cada configuración cada solución y permite evaluar cualquier
    \textit{schedule} sin volver a ejecutar los \textit{solvers}.
  \item Un conjunto de 20 características propias del JSSP, que se calculan en
    menos de un segundo y obtienen resultados comparables a los de
    \texttt{mzn2feat}.
  \item Una comparación controlada de estrategias de portafolio, que muestra
    que, con este portafolio y estas instancias, la selección de una
    configuración es mejor en
    promedio y el \textit{schedule} es más robusto frente a errores de
    selección.
  \item Un repositorio reproducible con las instancias, los modelos, el código,
    los resultados de todas las campañas y los programas que generan las tablas
    y figuras de este documento.
\end{itemize}

%------------------------------------------------------------------------------
\section{Organización del documento}
\label{sec:organizacion}
%------------------------------------------------------------------------------

El resto del documento se organiza en siete capítulos y cuatro anexos. El
Capítulo~\ref{cap:marco} presenta el marco de referencia sobre el JSSP, la
programación por restricciones, el problema de selección de algoritmos, los
portafolios de algoritmos, el uso del aprendizaje automático para predecir el
desempeño de los algoritmos y el estado del arte. El
Capítulo~\ref{cap:problema} define formalmente el problema y la metodología de
trabajo. El Capítulo~\ref{cap:caracterizacion} describe las instancias, la
extracción de características y la construcción del conjunto de datos. El
Capítulo~\ref{cap:portafolio} presenta el diseño y la implementación del
portafolio, incluidas las estrategias \textit{Winner Takes All} y de
\textit{schedule} secuencial. El Capítulo~\ref{cap:diseno} expone el diseño
experimental y el Capítulo~\ref{cap:resultados}, los resultados y su análisis.
Finalmente, el Capítulo~\ref{cap:conclusiones} recoge las conclusiones, las
limitaciones y el trabajo futuro. Los anexos contienen los resultados por
instancia, el catálogo de características, la configuración de los
experimentos y una guía de reproducción.
